% ============================================================
% Question Bank: ch08-07 Data Displays Extended
% Topic Title: Data Displays Extended
% CCSS: Supplementary
% Questions: 30 (18 MC + 7 SA + 5 Graphical)
% ============================================================

% --- Multiple Choice Questions (18) ---

\begin{practiceQuestion}{8-7-q01}{mc}
\begin{questionText}
A frequency table lists:
\end{questionText}
\choiceA{Only the largest and smallest values}
\choiceB{Categories and how often each occurs}
\choiceC{The mean, median, and mode}
\choiceD{Data arranged in a circle}
\correctAnswer{B}
\explanation{A frequency table shows each category (or interval) and the count of data values in it.}
\end{practiceQuestion}

\begin{practiceQuestion}{8-7-q02}{mc}
\begin{questionText}
Which data display is best for showing how data changes over time?
\end{questionText}
\choiceA{Circle graph}
\choiceB{Line graph}
\choiceC{Stem-and-leaf plot}
\choiceD{Box plot}
\correctAnswer{B}
\explanation{A line graph connects data points over time, making trends easy to see.}
\end{practiceQuestion}

\begin{practiceQuestion}{8-7-q03}{mc}
\begin{questionText}
A double bar graph is used to:
\end{questionText}
\choiceA{Show one data set over time}
\choiceB{Compare two data sets using side-by-side bars}
\choiceC{Show parts of a whole}
\choiceD{Display individual data points}
\correctAnswer{B}
\explanation{Double bar graphs place two sets of bars side by side for each category so you can compare them directly.}
\end{practiceQuestion}

\begin{practiceQuestion}{8-7-q04}{mc}
\begin{questionText}
A frequency table shows: Grade A: $10$, Grade B: $15$, Grade C: $8$, Grade D: $2$. What is the total number of students?
\end{questionText}
\choiceA{$25$}
\choiceB{$30$}
\choiceC{$35$}
\choiceD{$40$}
\correctAnswer{C}
\explanation{$10 + 15 + 8 + 2 = 35$.}
\end{practiceQuestion}

\begin{practiceQuestion}{8-7-q05}{mc}
\begin{questionText}
A line graph shows sales of $\$100$ in January, $\$150$ in February, $\$120$ in March. Between which months did sales decrease?
\end{questionText}
\choiceA{January to February}
\choiceB{February to March}
\choiceC{January to March}
\choiceD{Sales never decreased}
\correctAnswer{B}
\explanation{Sales went from $\$150$ in February to $\$120$ in March — a decrease of $\$30$.}
\end{practiceQuestion}

\begin{practiceQuestion}{8-7-q06}{mc}
\begin{questionText}
A frequency table has intervals $0$--$9$, $10$--$19$, $20$--$29$, $30$--$39$ with frequencies $3, 7, 5, 2$. Which interval has the most data values?
\end{questionText}
\choiceA{$0$--$9$}
\choiceB{$10$--$19$}
\choiceC{$20$--$29$}
\choiceD{$30$--$39$}
\correctAnswer{B}
\explanation{The interval $10$--$19$ has a frequency of $7$, which is the highest.}
\end{practiceQuestion}

\begin{practiceQuestion}{8-7-q07}{mc}
\begin{questionText}
In a double bar graph, Team X scored $30$, $45$, $50$ over three games and Team Y scored $35$, $40$, $55$. What was Team Y's total?
\end{questionText}
\choiceA{$120$}
\choiceB{$125$}
\choiceC{$130$}
\choiceD{$135$}
\correctAnswer{C}
\explanation{$35 + 40 + 55 = 130$.}
\end{practiceQuestion}

\begin{practiceQuestion}{8-7-q08}{mc}
\begin{questionText}
A line graph shows temperatures of $50°$, $55°$, $60°$, $58°$, $65°$ over five days. What is the overall trend?
\end{questionText}
\choiceA{Steady decrease}
\choiceB{Generally increasing}
\choiceC{No change}
\choiceD{Completely random}
\correctAnswer{B}
\explanation{Despite a small dip from $60°$ to $58°$, the overall trend from $50°$ to $65°$ is increasing.}
\end{practiceQuestion}

\begin{practiceQuestion}{8-7-q09}{mc}
\begin{questionText}
Which display is best for comparing exact frequencies of categorical data?
\end{questionText}
\choiceA{Line graph}
\choiceB{Frequency table}
\choiceC{Box plot}
\choiceD{Dot plot}
\correctAnswer{B}
\explanation{A frequency table lists exact counts for each category, which is ideal for comparing categorical data.}
\end{practiceQuestion}

\begin{practiceQuestion}{8-7-q10}{mc}
\begin{questionText}
A double bar graph compares boys and girls in a math competition. Boys scored $80, 85, 90$ and girls scored $75, 90, 95$ over three rounds. Who had the higher average?
\end{questionText}
\choiceA{Boys ($85$)}
\choiceB{Girls ($86.7$)}
\choiceC{They tied}
\choiceD{Cannot be determined}
\correctAnswer{B}
\explanation{Boys: $\frac{80+85+90}{3} = 85$. Girls: $\frac{75+90+95}{3} = \frac{260}{3} \approx 86.7$. Girls are slightly higher.}
\end{practiceQuestion}

\begin{practiceQuestion}{8-7-q11}{mc}
\begin{questionText}
A frequency table shows heights of students in $5$-inch intervals. The interval $60$--$64$ has frequency $12$ and $65$--$69$ has frequency $8$. What fraction of these $20$ students is in the $60$--$64$ interval?
\end{questionText}
\choiceA{$\frac{2}{5}$}
\choiceB{$\frac{3}{5}$}
\choiceC{$\frac{8}{20}$}
\choiceD{$\frac{12}{20}$}
\correctAnswer{B}
\explanation{$\frac{12}{20} = \frac{3}{5}$. (Note: choices B and D are equivalent; $\frac{3}{5}$ is simplified.)}
\end{practiceQuestion}

\begin{practiceQuestion}{8-7-q12}{mc}
\begin{questionText}
What is the main advantage of a line graph over a bar graph?
\end{questionText}
\choiceA{It can show more categories}
\choiceB{It clearly shows trends and changes over time}
\choiceC{It is easier to draw}
\choiceD{It uses colors}
\correctAnswer{B}
\explanation{Line graphs connect data points, making trends (increases, decreases, stability) easy to see.}
\end{practiceQuestion}

\begin{practiceQuestion}{8-7-q13}{mc}
\begin{questionText}
A line graph shows a company's profits: Q1 $\$50{,}000$, Q2 $\$60{,}000$, Q3 $\$55{,}000$, Q4 $\$70{,}000$. What was the increase from Q1 to Q4?
\end{questionText}
\choiceA{$\$10{,}000$}
\choiceB{$\$15{,}000$}
\choiceC{$\$20{,}000$}
\choiceD{$\$70{,}000$}
\correctAnswer{C}
\explanation{$\$70{,}000 - \$50{,}000 = \$20{,}000$.}
\end{practiceQuestion}

\begin{practiceQuestion}{8-7-q14}{mc}
\begin{questionText}
When would you choose a double bar graph instead of two separate bar graphs?
\end{questionText}
\choiceA{When the data sets have different categories}
\choiceB{When you want to compare two data sets for the same categories}
\choiceC{When you want to show change over time}
\choiceD{When you want to show parts of a whole}
\correctAnswer{B}
\explanation{A double bar graph is designed to compare two groups for the same categories by placing their bars side by side.}
\end{practiceQuestion}

\begin{practiceQuestion}{8-7-q15}{mc}
\begin{questionText}
A frequency table shows test scores:

$90$--$100$: $5$ students, $80$--$89$: $12$ students, $70$--$79$: $8$ students, $60$--$69$: $3$ students, Below $60$: $2$ students.

What percent scored $80$ or above?
\end{questionText}
\choiceA{$40\%$}
\choiceB{$56.7\%$}
\choiceC{$60\%$}
\choiceD{$17\%$}
\correctAnswer{B}
\explanation{Total: $5+12+8+3+2 = 30$. Scored $80+$: $5+12 = 17$. Percent: $\frac{17}{30} \approx 56.7\%$.}
\end{practiceQuestion}

\begin{practiceQuestion}{8-7-q16}{mc}
\begin{questionText}
A line graph shows a plant's height each week: Week 1: $2$ cm, Week 2: $5$ cm, Week 3: $9$ cm, Week 4: $14$ cm. Between which weeks did the plant grow the most?
\end{questionText}
\choiceA{Week 1 to Week 2}
\choiceB{Week 2 to Week 3}
\choiceC{Week 3 to Week 4}
\choiceD{Growth was the same each week}
\correctAnswer{C}
\explanation{Week 1--2: $3$ cm. Week 2--3: $4$ cm. Week 3--4: $5$ cm. The most growth was Week 3 to 4.}
\end{practiceQuestion}

\begin{practiceQuestion}{8-7-q17}{mc}
\begin{questionText}
A double bar graph shows that Store A sold $200$ items in June and $250$ in July. Store B sold $180$ in June and $300$ in July. Which store had a bigger increase?
\end{questionText}
\choiceA{Store A (increase of $50$)}
\choiceB{Store B (increase of $120$)}
\choiceC{They had the same increase}
\choiceD{Neither store increased}
\correctAnswer{B}
\explanation{Store A: $250 - 200 = 50$. Store B: $300 - 180 = 120$. Store B had a bigger increase.}
\end{practiceQuestion}

\begin{practiceQuestion}{8-7-q18}{mc}
\begin{questionText}
Which statement is true about frequency tables?
\end{questionText}
\choiceA{They only work for numerical data}
\choiceB{They can organize both categorical and numerical data}
\choiceC{They always require intervals}
\choiceD{They show data as a graph}
\correctAnswer{B}
\explanation{Frequency tables can list categories (favorite color) or numerical intervals (test scores $70$--$79$).}
\end{practiceQuestion}

% --- Short Answer Questions (7) ---

\begin{practiceQuestion}{8-7-q19}{sa}
\begin{questionText}
A frequency table shows: Soccer $15$, Basketball $20$, Tennis $10$, Swimming $5$. What is the total and what percent play basketball?
\end{questionText}
\correctAnswer{Total $= 50$; Basketball $= 40\%$}
\explanation{$15+20+10+5 = 50$. Basketball: $\frac{20}{50} = 40\%$.}
\end{practiceQuestion}

\begin{practiceQuestion}{8-7-q20}{sa}
\begin{questionText}
A line graph shows monthly rainfall: Jan $2$ in., Feb $3$ in., Mar $5$ in., Apr $4$ in. Find the total rainfall and the month with the most.
\end{questionText}
\correctAnswer{Total $= 14$ in.; March had the most ($5$ in.)}
\explanation{$2+3+5+4 = 14$ inches. March has the highest value at $5$ in.}
\end{practiceQuestion}

\begin{practiceQuestion}{8-7-q21}{sa}
\begin{questionText}
In a double bar graph, Team A scores: $50, 60, 70$. Team B scores: $45, 65, 80$. Find the mean for each team.
\end{questionText}
\correctAnswer{Team A mean $= 60$; Team B mean $\approx 63.3$}
\explanation{Team A: $\frac{50+60+70}{3} = 60$. Team B: $\frac{45+65+80}{3} = \frac{190}{3} \approx 63.3$.}
\end{practiceQuestion}

\begin{practiceQuestion}{8-7-q22}{sa}
\begin{questionText}
A line graph shows temperatures: Mon $72°$, Tue $68°$, Wed $70°$, Thu $75°$, Fri $73°$. What is the range of temperatures?
\end{questionText}
\correctAnswer{$7°$}
\explanation{Range $= 75 - 68 = 7°$.}
\end{practiceQuestion}

\begin{practiceQuestion}{8-7-q23}{sa}
\begin{questionText}
A frequency table shows:

$1$--$5$: $4$, $6$--$10$: $9$, $11$--$15$: $6$, $16$--$20$: $1$. What percent of data values are in the $6$--$10$ interval?
\end{questionText}
\correctAnswer{$45\%$}
\explanation{Total: $4+9+6+1 = 20$. Percent: $\frac{9}{20} = 45\%$.}
\end{practiceQuestion}

\begin{practiceQuestion}{8-7-q24}{sa}
\begin{questionText}
Give one advantage and one disadvantage of using a frequency table compared to a bar graph.
\end{questionText}
\correctAnswer{Advantage: shows exact counts; Disadvantage: harder to see trends or compare visually}
\explanation{Tables give precise numbers but graphs make patterns and comparisons easier to spots at a glance.}
\end{practiceQuestion}

\begin{practiceQuestion}{8-7-q25}{sa}
\begin{questionText}
A double bar graph shows monthly book sales: Store A sold $120$ in March and $150$ in April. Store B sold $100$ in March and $180$ in April. Which store had a greater percent increase from March to April?
\end{questionText}
\correctAnswer{Store B ($80\%$ increase vs.\ $25\%$)}
\explanation{Store A: $\frac{150-120}{120} = \frac{30}{120} = 25\%$. Store B: $\frac{180-100}{100} = \frac{80}{100} = 80\%$. Store B had a greater percent increase.}
\end{practiceQuestion}

% --- Graphical Questions (3 GMC + 2 GSA) ---

\begin{practiceQuestion}{8-7-q26}{gmc}
\begin{questionText}
The frequency table shows the number of books students read last month.

\begin{center}
\begin{tabular}{|c|c|}
\hline
\textbf{Books Read} & \textbf{Frequency} \\
\hline
$0$--$2$ & $6$ \\
\hline
$3$--$5$ & $14$ \\
\hline
$6$--$8$ & $8$ \\
\hline
$9$--$11$ & $2$ \\
\hline
\end{tabular}
\end{center}

What percent of students read $3$ or more books?
\end{questionText}
\choiceA{$46.7\%$}
\choiceB{$73.3\%$}
\choiceC{$80\%$}
\choiceD{$20\%$}
\correctAnswer{C}
\explanation{Total: $6+14+8+2 = 30$. Read $3+$: $14+8+2 = 24$. Percent: $\frac{24}{30} = 80\%$.}
\end{practiceQuestion}

\begin{practiceQuestion}{8-7-q27}{gmc}
\begin{questionText}
The double bar graph shows the number of goals scored by two soccer teams over four games.

\begin{center}
\begin{tikzpicture}[scale=0.55]
  \draw[thick,->] (0,0) -- (0,6) node[above]{\small Goals};
  \draw[thick,->] (0,0) -- (10,0);
  % Game 1
  \fill[funBlue!70] (0.5,0) rectangle (1.2,3);
  \fill[funGreen!70] (1.3,0) rectangle (2.0,2);
  \node[below,font=\tiny] at (1.25,0) {G1};
  % Game 2
  \fill[funBlue!70] (2.5,0) rectangle (3.2,4);
  \fill[funGreen!70] (3.3,0) rectangle (4.0,5);
  \node[below,font=\tiny] at (3.25,0) {G2};
  % Game 3
  \fill[funBlue!70] (4.5,0) rectangle (5.2,2);
  \fill[funGreen!70] (5.3,0) rectangle (6.0,3);
  \node[below,font=\tiny] at (5.25,0) {G3};
  % Game 4
  \fill[funBlue!70] (6.5,0) rectangle (7.2,5);
  \fill[funGreen!70] (7.3,0) rectangle (8.0,4);
  \node[below,font=\tiny] at (7.25,0) {G4};
  \foreach \y in {1,...,5} { \draw (-0.1,\y) -- (0.1,\y); \node[left,font=\tiny] at (0,\y) {$\y$}; }
  % legend
  \fill[funBlue!70] (8.5,5) rectangle (9,5.3);
  \node[right,font=\tiny] at (9,5.15) {Team A};
  \fill[funGreen!70] (8.5,4.3) rectangle (9,4.6);
  \node[right,font=\tiny] at (9,4.45) {Team B};
\end{tikzpicture}
\end{center}

Which team scored more total goals?
\end{questionText}
\choiceA{Team A with $14$ goals}
\choiceB{Team B with $14$ goals}
\choiceC{They tied with $14$ goals each}
\choiceD{Team A with $12$ goals}
\correctAnswer{C}
\explanation{Team A: $3+4+2+5 = 14$. Team B: $2+5+3+4 = 14$. They tied.}
\end{practiceQuestion}

\begin{practiceQuestion}{8-7-q28}{gmc}
\begin{questionText}
The line graph shows a student's weekly quiz scores over $5$ weeks.

\begin{center}
\begin{tikzpicture}[scale=0.55]
  \draw[thick,->] (0,0) -- (0,6) node[above]{\small Score};
  \draw[thick,->] (0,0) -- (7,0);
  \foreach \x/\lab in {1/W1,2/W2,3/W3,4/W4,5/W5} {
    \draw (\x,0.1) -- (\x,-0.1);
    \node[below,font=\tiny] at (\x,0) {\lab};
  }
  \foreach \y in {1,...,5} { \draw (-0.1,\y) -- (0.1,\y); \node[left,font=\tiny] at (0,\y) {$\y\times 20$}; }
  % Scores: 60, 70, 65, 80, 90
  \fill[funBlue] (1,3) circle (0.1); % 60
  \fill[funBlue] (2,3.5) circle (0.1); % 70
  \fill[funBlue] (3,3.25) circle (0.1); % 65
  \fill[funBlue] (4,4) circle (0.1); % 80
  \fill[funBlue] (5,4.5) circle (0.1); % 90
  \draw[thick,funBlue] (1,3) -- (2,3.5) -- (3,3.25) -- (4,4) -- (5,4.5);
  \node[above,font=\tiny] at (1,3) {$60$};
  \node[above,font=\tiny] at (2,3.5) {$70$};
  \node[above right,font=\tiny] at (3,3.25) {$65$};
  \node[above,font=\tiny] at (4,4) {$80$};
  \node[above,font=\tiny] at (5,4.5) {$90$};
\end{tikzpicture}
\end{center}

Between which two consecutive weeks was the greatest improvement?
\end{questionText}
\choiceA{Week 1 to Week 2}
\choiceB{Week 2 to Week 3}
\choiceC{Week 3 to Week 4}
\choiceD{Week 4 to Week 5}
\correctAnswer{C}
\explanation{W1--W2: $+10$. W2--W3: $-5$. W3--W4: $+15$. W4--W5: $+10$. The biggest improvement is $+15$ from Week 3 to Week 4.}
\end{practiceQuestion}

\begin{practiceQuestion}{8-7-q29}{gsa}
\begin{questionText}
The double bar graph compares test scores for two classes over three tests.

\begin{center}
\begin{tikzpicture}[scale=0.5]
  \draw[thick,->] (0,0) -- (0,6.5) node[above]{\small Average Score};
  \draw[thick,->] (0,0) -- (8.5,0);
  % Test 1
  \fill[funBlue!70] (0.5,0) rectangle (1.5,4);
  \fill[funOrange!70] (1.6,0) rectangle (2.6,3.5);
  \node[below,font=\tiny] at (1.55,0) {Test 1};
  % Test 2
  \fill[funBlue!70] (3.2,0) rectangle (4.2,4.5);
  \fill[funOrange!70] (4.3,0) rectangle (5.3,5);
  \node[below,font=\tiny] at (4.25,0) {Test 2};
  % Test 3
  \fill[funBlue!70] (5.9,0) rectangle (6.9,5);
  \fill[funOrange!70] (7.0,0) rectangle (8.0,4);
  \node[below,font=\tiny] at (6.95,0) {Test 3};
  \foreach \y in {1,...,6} { \draw (-0.1,\y) -- (0.1,\y); \node[left,font=\tiny] at (0,\y) {$\y\times 15$}; }
  \node[above,font=\tiny] at (1,4) {$60$};
  \node[above,font=\tiny] at (2.1,3.5) {$52$};
  \node[above,font=\tiny] at (3.7,4.5) {$68$};
  \node[above,font=\tiny] at (4.8,5) {$75$};
  \node[above,font=\tiny] at (6.4,5) {$75$};
  \node[above,font=\tiny] at (7.5,4) {$60$};
\end{tikzpicture}
\end{center}

Class A (blue): $60, 68, 75$. Class B (orange): $52, 75, 60$. Find the mean for each class. Which class improved more consistently?
\end{questionText}
\correctAnswer{Class A mean $\approx 67.7$; Class B mean $= 62.3$. Class A improved consistently (scores went up each test).}
\explanation{Class A: $\frac{60+68+75}{3} = \frac{203}{3} \approx 67.7$. Class B: $\frac{52+75+60}{3} = \frac{187}{3} \approx 62.3$. Class A's scores increased each test; Class B went up then down.}
\end{practiceQuestion}

\begin{practiceQuestion}{8-7-q30}{gsa}
\begin{questionText}
The line graph shows the number of visitors to a library each day for one week.

\begin{center}
\begin{tikzpicture}[scale=0.5]
  \draw[thick,->] (0,0) -- (0,7) node[above]{\small Visitors};
  \draw[thick,->] (0,0) -- (9,0);
  \foreach \x/\lab in {1/Mon,2/Tue,3/Wed,4/Thu,5/Fri,6/Sat,7/Sun} {
    \draw (\x,0.1) -- (\x,-0.1);
    \node[below,font=\tiny,rotate=45,anchor=north east] at (\x,0) {\lab};
  }
  \foreach \y in {1,...,6} { \draw (-0.1,\y) -- (0.1,\y); \node[left,font=\tiny] at (0,\y) {$\y\times 10$}; }
  % Data: Mon 40, Tue 35, Wed 50, Thu 45, Fri 60, Sat 55, Sun 30
  \fill[funBlue] (1,4) circle (0.12);
  \fill[funBlue] (2,3.5) circle (0.12);
  \fill[funBlue] (3,5) circle (0.12);
  \fill[funBlue] (4,4.5) circle (0.12);
  \fill[funBlue] (5,6) circle (0.12);
  \fill[funBlue] (6,5.5) circle (0.12);
  \fill[funBlue] (7,3) circle (0.12);
  \draw[thick,funBlue] (1,4) -- (2,3.5) -- (3,5) -- (4,4.5) -- (5,6) -- (6,5.5) -- (7,3);
  \node[above,font=\tiny] at (1,4) {$40$};
  \node[above,font=\tiny] at (2,3.5) {$35$};
  \node[above,font=\tiny] at (3,5) {$50$};
  \node[above,font=\tiny] at (4,4.5) {$45$};
  \node[above,font=\tiny] at (5,6) {$60$};
  \node[above,font=\tiny] at (6,5.5) {$55$};
  \node[above,font=\tiny] at (7,3) {$30$};
\end{tikzpicture}
\end{center}

Find the mean number of daily visitors and the range. On which day was the library busiest?
\end{questionText}
\correctAnswer{Mean $= 45$ visitors; Range $= 30$; Friday was busiest ($60$ visitors)}
\explanation{Sum: $40+35+50+45+60+55+30 = 315$. Mean: $\frac{315}{7} = 45$. Range: $60 - 30 = 30$. Friday had the most visitors ($60$).}
\end{practiceQuestion}
